%% A short paper in the shape of the acmart sigconf sample (SIGGRAPH conference track).
%% Written for the renderer's golden tests; every construct here has a test that pins it.
\documentclass[sigconf]{acmart}
\usepackage{booktabs}
\usepackage{tikz}

\citestyle{acmauthoryear}

\newcommand{\ours}{LumenField}
\newcommand{\vect}[1]{\mathbf{#1}}
\newtheorem{theorem}{Theorem}

\setcopyright{acmlicensed}
\copyrightyear{2026}
\acmYear{2026}
\acmDOI{10.1145/0000000.0000000}
\acmConference[SIGGRAPH '26]{Special Interest Group on Computer Graphics and Interactive Techniques Conference}{July 2026}{Vancouver, BC, Canada}
\acmISBN{978-1-4503-XXXX-X/26/07}

\begin{document}

\title[\ours{}: Radiance Fields]{\ours{}: Radiance Fields from a Single Photograph}

\author{Ada Example}
\orcid{0000-0002-1825-0097}
\email{ada@example.edu}
\affiliation{%
  \institution{Institute of Light Transport}
  \city{Vancouver}
  \country{Canada}
}

\author{Fran\c{c}ois M\"uller}
\authornote{Both authors contributed equally to this research.}
\email{francois@example.org}
\affiliation{%
  \institution{Example Graphics Lab}
  \city{Z\"urich}
  \country{Switzerland}
}

\begin{abstract}
We present \ours{}, a method that recovers a radiance field from one photograph
in under a second. Prior work needs dozens of views~\cite{mildenhall2020nerf};
ours needs one, and it runs on a laptop.
\end{abstract}

\begin{CCSXML}
<ccs2012>
   <concept>
       <concept_id>10010147.10010371.10010372</concept_id>
       <concept_desc>Computing methodologies~Rendering</concept_desc>
       <concept_significance>500</concept_significance>
   </concept>
</ccs2012>
\end{CCSXML}

\ccsdesc[500]{Computing methodologies~Rendering}
\ccsdesc[300]{Computing methodologies~Neural networks}

\keywords{radiance fields, single-view reconstruction, real-time rendering}

\begin{teaserfigure}
  \includegraphics[width=\textwidth]{figures/teaser.png}
  \caption{\ours{} turns one photograph (left) into a navigable scene (right).}
  \Description{A photograph of a room next to three rendered viewpoints of the same room.}
  \label{fig:teaser}
\end{teaserfigure}

\maketitle

\section{Introduction}
\label{sec:intro}

Radiance fields changed how we think about scene capture~\cite{mildenhall2020nerf, kerbl2023gaussians}.
They also made capture slow: a good field needs many photographs and minutes of
optimisation. \ours{} removes both costs. As Figure~\ref{fig:teaser} shows, a single
photograph is enough, and \citet{kerbl2023gaussians} report similar quality only with
hundreds of views.

Our contributions are:
\begin{itemize}
  \item a feed-forward predictor for the field parameters (Section~\ref{sec:method});
  \item a training objective that needs no ground-truth geometry;
  \item an evaluation on three public datasets (Table~\ref{tab:results}).
\end{itemize}

\section{Method}
\label{sec:method}

Given an image $I$ we predict a density $\sigma(\vect{x})$ and a colour
$c(\vect{x}, \vect{d})$ for every point $\vect{x}$ and direction $\vect{d}$. The
rendered colour of a ray is
\begin{equation}
  C(\mathbf{r}) = \int_{t_n}^{t_f} T(t)\,\sigma(\mathbf{r}(t))\,c(\mathbf{r}(t), \mathbf{d})\,dt,
  \label{eq:volume}
\end{equation}
where $T(t)$ is the accumulated transmittance. Equation~\eqref{eq:volume} is
differentiable, so the predictor trains end to end.\footnote{We use the quadrature
of \citet{mildenhall2020nerf} with 64 samples per ray.}

\begin{theorem}[Consistency]
  If the predictor is Lipschitz then the rendered views are consistent under small
  camera motion.
\end{theorem}

\subsection{Predictor}

The predictor is a transformer over image patches. We write its layers as
\begin{align}
  h_0 &= \mathrm{Embed}(I), \\
  h_{\ell+1} &= h_\ell + \mathrm{Attn}(h_\ell) \label{eq:layer}
\end{align}
and read the field parameters from $h_L$ (see Equation~\eqref{eq:layer}).

\section{Results}

\begin{table}
  \caption{PSNR on three datasets. Higher is better.}
  \label{tab:results}
  \begin{tabular}{lccc}
    \toprule
    Method & LLFF & DTU & Mip-360 \\
    \midrule
    NeRF & 26.5 & 27.1 & 24.9 \\
    3DGS & 27.2 & 28.0 & 27.4 \\
    \ours{} & \textbf{27.9} & \textbf{28.6} & 27.1 \\
    \bottomrule
  \end{tabular}
\end{table}

Table~\ref{tab:results} reports PSNR; the live numbers are kept in the dynamic table
below, which updates without a new version.

\begin{table}
  \caption{Ablations, updated as runs finish.}
  \label{tab:ablation}
  \derivetable{ablation}
\end{table}

\begin{figure}
  \derivefigure[width=0.8\linewidth]{qualitative}
  \caption{Qualitative results on held-out scenes.}
  \label{fig:qual}
\end{figure}

We also compare against \emph{fast} baselines --- see \textbf{Figure~\ref{fig:qual}} and
Figure~\ref{fig:teaser}. ``Quoted'' text and a non-breaking~space survive.

\begin{acks}
We thank the anonymous reviewers. This work was supported by grant~1234.
\end{acks}

\bibliographystyle{ACM-Reference-Format}
\bibliography{refs}

\appendix

\section{Implementation Details}

We train for 100k iterations with Adam.

\end{document}
